\par\Needspace{14\baselineskip}\medskip
\begin{center}\hypertarget{native-230-article167}{}\pdfbookmark[2]{167.}{article-167}167.\end{center}
Before explaining the solution, we give a list of the principal symbols. We shall use
\par\noindent
\begin{tabular}{@{}ll@{}}
$t,\ t',\ t'',\ t'''$ & the times of the four observations\\
$a,\ a',\ a'',\ a'''$ & the geocentric longitudes of the celestial body\\
$\beta,\ \beta',\ \beta'',\ \beta'''$ & its latitudes\\
$r,\ r',\ r'',\ r'''$ & the distances from the Sun\\
$\rho,\ \rho',\ \rho'',\ \rho'''$ & the distances from the Earth\\
$l,\ l',\ l'',\ l'''$ & the heliocentric longitudes of the Earth\\
$B,\ B',\ B'',\ B'''$ & the heliocentric latitudes of the Earth\\
$R,\ R',\ R'',\ R'''$ & the distances of the Earth from the Sun
\end{tabular}
\par\noindent\SourceLeaf{231}
$(n01),\ (n12),\ (n23),\ (n02),\ (n13)$ for twice the areas of the triangles formed respectively
by the Sun and the first and second,
second and third, third and fourth, first and third, and second and fourth positions of the celestial body
\par\noindent
$\dfrac{1}{(\eta01)}\mathpunct{\raisebox{.45ex}{,}}\quad
\dfrac{1}{(\eta12)}\mathpunct{\raisebox{.45ex}{,}}\quad
\dfrac{1}{(\eta23)}$ for the quotients obtained by dividing the areas
\newline
$\frac12(n01),\ \frac12(n12),\ \frac12(n23)$
by the areas of the corresponding sectors
\[
P'=\frac{(n12)}{(n01)}\mathpunct{\raisebox{.45ex}{,}}\qquad
P''=\frac{(n12)}{(n23)}
\]
\[
\begin{aligned}
Q'&=\left(\frac{(n01)+(n12)}{(n02)}-1\right)r'^3,\\[.5em]
Q''&=\left(\frac{(n12)+(n23)}{(n13)}-1\right)r''^3
\end{aligned}
\]
\noindent $v,\ v',\ v'',\ v'''$ for the longitudes of the celestial body in its orbit, reckoned from an arbitrary point.
\par
Finally, for the second and third observations we shall denote the heliocentric positions of the Earth on
the celestial sphere by $A'$, $A''$, the geocentric positions of the celestial body
by $B'$, $B''$, and its heliocentric positions by $C'$, $C''$.
\par
With these conventions, our first task, as in the problem of the preceding section
(Article 136), will be to determine the orientation of the great circles $A'C'B'$,
$A''C''B''$, whose inclinations to the ecliptic we denote by $\gamma'$, $\gamma''$. In the same
calculation we also determine the arcs $A'B'=\delta'$, $A''B''=\delta''$. From this
it clearly follows that
\[
\begin{aligned}
r'&=\surd(\rho'\rho'+2\rho'R'\cos\delta'+R'R')\\
r''&=\surd(\rho''\rho''+2\rho''R''\cos\delta''+R''R''),
\end{aligned}
\]
\noindent or, setting $\rho'+R'\cos\delta'=x'$, $\rho''+R''\cos\delta''=x''$,
$R'\sin\delta'=a'$, $R''\sin\delta''=a''$, we have
\[
\begin{aligned}
r'&=\surd(x'x'+a'a')\\
r''&=\surd(x''x''+a''a'').
\end{aligned}
\]
